\documentclass[letterpaper,11pt]{article}
\usepackage{graphicx}
\usepackage{geometry}
\linespread{1.2}                        
\setlength{\oddsidemargin}{0.0in}            
\setlength{\textwidth}{6.5in}    
\usepackage[utf8]{inputenc} % Asegúrate de usar utf8
\usepackage[T1]{fontenc} % Asegúrate de usar T1 font encoding      
\setlength{\topmargin}{-.6in}            
\setlength{\textheight}{9in}
\usepackage{latexsym} 
\usepackage{amssymb} 
\usepackage{amsmath}
\usepackage{amsxtra}
\usepackage[mathcal]{eucal} 
\usepackage{enumerate}
\usepackage{hyperref}
\usepackage{rotating}
\usepackage{caption}
\usepackage{lscape}
\usepackage{paralist} 
\usepackage{indentfirst}
\usepackage[dvipsnames,svgnames]{xcolor} 
\usepackage{setspace}
\usepackage{multicol} 
\usepackage{float}

\definecolor{codegreen}{rgb}{0,0.6,0}
\definecolor{codegray}{rgb}{0.5,0.5,0.5}
\definecolor{codepurple}{cmyk}{0.65, 1, 0, 0.35}
\definecolor{backcolour}{rgb}{0.95,0.95,0.92}

\usepackage{listings}
\usepackage{xcolor}

\lstset{
    language=R, % lenguaje de programación
    basicstyle=\small\ttfamily, % estilo básico
    commentstyle=\color{gray}, % estilo para los comentarios
    keywordstyle=\color{blue}, % estilo para las palabras clave
    stringstyle=\color{black}, % estilo para las cadenas de texto
    showstringspaces=false, % no mostrar espacios en cadenas de texto
    tabsize=4, % tamaño de la tabulación
    breaklines=true, % permitir el salto de línea automático
    postbreak=\mbox{\textcolor{red}{$\hookrightarrow$}\space}, % símbolo al final de una línea partida
    numbers=left, % mostrar números de línea
    numberstyle=\tiny\color{gray}, % estilo para los números de línea
    frame=single, % mostrar un borde alrededor del código
    backgroundcolor=\color{gray!10}, % color de fondo
}

\begin{document}

\begin{center}
    {\Large  Taller 1: Econometría }

Profesor:    Pedro Comte\\
Ayudantes: Jociney Honorio y Fabian Cisternas

\end{center}

\textbf{Instrucciones:} 

\begin{itemize}
    \item Usted dispone de 60 minutos  para intentar los 100 puntos del taller.
    \item Cuide la presentación y redacción de sus respuestas.
    \item Puede utilizar su computador y los apuntes tanto de clase como de ayudantía.
    \item Debe enviar un archivo en formato PDF y un R script (extensión .R)
\end{itemize}

\section{Datos Aleatorios (30 puntos)}
\begin{enumerate}
    \item Genere el dataframe "datos" utilizando el código en Anexos. (6 pts.)

Según la letra por la que comience su apellido, cambie (o mantenga) la semilla:

\begin{table}[htbp]
\centering
\begin{tabular}{|c|c|c|c|c|}
\hline
\textbf{A-E} & \textbf{F-J} & \textbf{K-O} & \textbf{P-T} & \textbf{U-Z} \\
\hline
123 & 456 & 789 & 101112 & 131415 \\
\hline
\end{tabular}
% \caption{División del alfabeto en 5 segmentos con semillas asignadas}
\label{tabla:division-alfabeto}
\end{table}

El set de datos generado explica las notas de 50 alumnos de la USM:

\begin{itemize}
    \item \textbf{Horas de sueño}: Horas de sueño promedio que cada alumno ha tenido en el último mes.
    \item \textbf{Profesor particular}: Indica si cada alumno ha recibido ayuda externa de un profesor particular o no.
    \item \textbf{Promedio notas semestre pasado}: Promedio de notas que cada alumno obtuvo en el semestre pasado.
    \item \textbf{Tiempo de estudio}: Tiempo de estudio en horas que cada alumno dedica.
    \item \textbf{Asistencia a clases}: Porcentaje de asistencia a clases de cada alumno.
    \item \textbf{Nivel socioeconomico}: Nivel socioeconómico de cada alumno, con valores del 1 al 5.
\end{itemize}

A partir del dataset generado:
    \item Redondee las notas a 1 decimal. (2 pts.)
\begin{lstlisting}
datos$notas <- round(datos$notas, 1)
\end{lstlisting}


    \item Calcule manualmente los beta estimadores a través de la matriz X y el vector de la variable dependiente (Sin usar funciones en R). (8 pts.)
\begin{lstlisting}
# Calcular X
X <- cbind(rep(1, nrow(datos)), datos$hrs_sueno, datos$profesor_part, datos$media_sem_pasado, datos$tiempo_est, datos$asistencia, datos$nivel_socioec)

# Calcular X^T
Xt <- t(X)

# Calcular X^T X
XtX <- t(X) %*% X

# Calcular X^T y
Xty <- t(X) %*% datos$notas

# Calcular los estimadores beta
beta_hat <- solve(XtX) %*% Xty
\end{lstlisting}



    \item Genera un \textit{resumen} a partir de un modelo de regresión múltiple, esta vez utilizando funciones en R e interprete los \textit{betas} del modelo. (8 pts.)
    \begin{lstlisting}
# Generar modelo utilizando función LM
nombres_vars <- c("hrs_sueno", "profesor_part", "media_sem_pasado", "tiempo_est", "asistencia", "nivel_socioec")
modelo <- lm(datos$notas ~ ., data = datos[, nombres_vars])
summary(modelo)

# Ver el resumen del modelo
summary(modelo)
\end{lstlisting}

        \item \textquestiondown Como se relacionan los \textit{betas} obtenidos con las ponderaciones para construir las notas en el codigo anexo? (6 pts.)
        


    
\end{enumerate}

\section{Wooldridge  (70 puntos)}

\subsection{Base Wage (30 puntos)}

Utilice los datos en el archivo WAGE1 de Wooldridge para determinar los estimadores asociados al  siguiente modelo de regresión poblacional.

$wage=\beta_0+ \beta_1educ +\beta_2exper+ \beta_3tenure+u$


Donde wage es salario, educ son los años de escolaridad, 
exper son los años de experiencia laboral, y tenure los años en su actual trabajo.

\begin{enumerate}


\begin{lstlisting}
# Wooldridge
library(wooldridge)
data("wage1")
wage1 <- na.omit(wage1)
\end{lstlisting}
\item Determine los estimadores mediante MCO, e interprete sus resultados. (8 puntos)
\begin{lstlisting}
# modelo de regresión y summary
modelo <- lm(wage ~ educ + exper+tenure, data=wage1)
summary(modelo)
\end{lstlisting}


\item \textquestiondown Los estimadores tienen los signos que usted esperaba? (10 puntos)


\item Haga un nuevo modelo de regresión lineal, pero utilizando solo las variables independientes \textit{educ} y \textit{tenure}. Comparar $R^2$ y $R^2_{ajustado}$ entre el modelo completo y el nuevo modelo \textquestiondown Qué se puede concluir?. (12 puntos)
\begin{lstlisting}
# modelo de regresión y summary
modelo <- lm(wage ~ educ +tenure, data=wage1)
summary(modelo)
\end{lstlisting}



\end{enumerate}

\subsection{Base Attend (40 puntos)}

\begin{enumerate}
    \item Utilizando R, carge la base de datos "\textit{attend}" desde Wooldridge. (4 puntos)

\begin{lstlisting}
# Wooldridge
library(wooldridge)
data("attend")
attend <- na.omit(attend)
\end{lstlisting}

    \item Sub seleccione las siguientes variables y almacénelas en un nuevo set de datos (4 puntos):

\begin{itemize}
\item \textbf{attend}: Clases asistidas de un total de 32.
\item \textbf{termGPA}: Promedio de calificaciones durante el período.
\item \textbf{priGPA}: Promedio de calificaciones acumulado antes de este período.
\item \textbf{ACT}: puntaje en el examen ACT.
\item \textbf{final}: puntaje del examen final.
\item \textbf{hwrte}: porcentaje de tareas entregadas.
\item \textbf{frosh}: =1 si es estudiante de primer año, 0 si no.
\item \textbf{soph}:=1 si es estudiante de segundo año, 0 si no.
\end{itemize}
\begin{lstlisting}
var_selec <- c("attend","termGPA","priGPA","ACT","final","hwrte","frosh","soph")
attend_new <- subset(attend, select=var_selec)
\end{lstlisting}


    \item Genere un modelo de regresión lineal múltiple que contenga las variables anteriores, donde la variable dependiente es el puntaje final en el examen, y un resumen del modelo utilizando \textit{summary}. (10 puntos)
 \begin{lstlisting}
modelo <- lm(final ~ . , data=attend_new)
summary(modelo)
\end{lstlisting}   

    \item Interprete el resultado de la bondad de ajuste y su versión ajustada. (6 puntos)

    \item Genere un nuevo modelo pero excluyendo las variables no estadísticamente significativas (\textit{p-value > 0.05}). (10 puntos)
 \begin{lstlisting}
\begin{lstlisting}
var_selec <- c("termGPA","ACT","final","soph")
attend_new2 <- subset(attend, select=var_selec)
\end{lstlisting}


    \item Compare su nuevo modelo con su anterior modelo en términos bondad de ajuste en su versión ajustada. (6 puntos)
 \begin{lstlisting}
modelo2 <- lm(final ~. , data=attend_new2)
summary(modelo2)

\end{lstlisting}

\end{enumerate}
\newpage
\section{Anexos: Código}
\begin{verbatim}
# Establecer la semilla para reproducibilidad
set.seed(123)

# Crear la base de datos ficticia con 50 alumnos
datos <- data.frame(
  hrs_sueno = round(runif(50, min = 5, max = 10), 1),
  profesor_part = sample(c(0, 1), 50, replace = TRUE),
  media_sem_pasado = round(runif(50, min = 60, max = 100), 1),
  tiempo_est = round(runif(50, min = 1, max = 8), 1),
  asistencia = round(runif(50, min = 60, max = 100), 1),
  nivel_socioec = sample(1:5, 50, replace = TRUE)
)

# Calcular las notas con ajustes en la ponderación
datos$notas <- 30 + # Base más baja para empezar desde un mínimo más realista
  datos$hrs_sueno * 1.5 + # Reducir el impacto del sueño
  datos$profesor_part * 3 + # Menos puntos adicionales por participación del profesor
  datos$media_sem_pasado * 0.2 + # Menor peso de la media semestral pasada
  datos$tiempo_est * 2 + # Ajuste en tiempo de estudio
  datos$asistencia * 0.15 + # Menor impacto de la asistencia
  datos$nivel_socioec * 2 + # Menos impacto del nivel socioeconómico
  rnorm(50, mean = 0, sd = 5) # Menor desviación estándar para reducir variabilidad

# Asegurar que las notas estén entre 20 y 100
datos$notas <- pmax(pmin(datos$notas, 100), 20)
\end{verbatim}
\end{document}






